\begin{enumerate}
\item 
	\begin{enumerate}
	\item  	\textbf{Affirmation 1  est Fausse. } 
	
	L'équation différentielle \quad $\left(E_{1}\right): \quad y'+2 y=0 \iff y'=-2y$ est de la forme $y'=ay$ 
	
donc les  solutions sont les fonctions de la forme $x\mapsto C\e^{ax}$ avec $C\in\R$. 

Donc, il existe $C\in \R$ tel que $f(x)= C \e^{-2x}$ .

Or $f(0)=1$ donc $1=C\e^{-2\times 0} \iff C=1$ .

donc $f$ est définie sur $\R$ par $f(x)=\e^{-2x}$. 

$f(ln 2) = \e^{-2\ln 2} = \e^{ -ln 2^2}= \e^{ln \left(\frac{1}{4}\right)}= \dfrac{1}{4} \neq \dfrac{1}{2}$

Donc : le point A de coordonnées $\left(\ln 2 ; \dfrac{1}{2}\right)$ n'appartient pas à la courbe représentative de la fonction $f$.



	\item	\textbf{Affirmation 2  est  Fausse. } 
	
Pour étudier la position relative de la courbe représentative de la fonction $f$ et de ses tangentes, étudions la convexité de la fonction. 

Pour tout réel $x\in \left[\dfrac{\pi}{2} ; \pi\right]$  : 

$f'(x) = 1 + \sin(x)$ 

et donc   $f''(x)= \cos(x)$.

Or sur l'intervalle $\left[\dfrac{\pi}{2} ; \pi\right]$ , $cos(x) \leq 0$ donc la fonction est concave. 

Elle est donc située en-dessous de toutes ses tangentes. 





	\item 	\textbf{Affirmation 3  est  Vraie. }
	
	
	Pour tout réel  $x\in ] 0 ;+\infty[$ : 
	
	$\begin{aligned}
		-1 \leqslant \sin(x) \leqslant 1 &\implies  -\dfrac{1}{x} \leqslant  \dfrac{\sin (x)}{x} \leqslant \dfrac{1}{x}  \quad \text{car } x>0 \\
	&	\implies  1 -\dfrac{1}{x} \leqslant  1 + \dfrac{\sin (x)}{x} \leqslant 1 +\dfrac{1}{x}  \\
	&	\implies  1 -\dfrac{1}{x} \leqslant  f(x)  \leqslant 1 +\dfrac{1}{x}  \\
	\end{aligned}$
	
	Or $\lim\limits_{x\to +\infty} \dfrac{1}{x} =0 $ donc, par somme,  $\lim\limits_{x\to +\infty }1 -\dfrac{1}{x} =1 $ et  $\lim\limits_{x\to +\infty }1 +\dfrac{1}{x} =1 $. 
	
	D'après le théorème des gendarmes, on a donc : $\lim\limits_{x\to +\infty }f(x) =1 $. 
	
	C'est à dire, la droite d'équation $y=1$ est une asymptote horizontale à la courbe représentative de la fonction $f$ au voisinage de $+\infty$.
	\end{enumerate}



	\item 	\textbf{Affirmation 4  est Vraie.} 
	
	
	 La fonction $f$ est dérivable sur $\R$ et pour tout réel $x$, $f'(x) = 2 -3 \e^{3x}$.
	
	Donc, pour tout réel $x$ : 
	
	$\begin{aligned} f'(x) -3f(x) &=  2 -3 \e^{3x} - 3 \left(2x-\e^{3x}\right) \\
		 &= 2 -3 \e^{3x} - 6x +3 \e^{3x} \\
		 &= 2  - 6x 
		 \end{aligned}$ 
		 
		 Donc $f$ est solution de l'équation $(E_{2})$.
	
\end{enumerate}
